\documentclass[12pt, a4paper]{article}

\usepackage[utf8]{inputenc}
\usepackage[T2A]{fontenc}
\usepackage[russian]{babel}
\usepackage{amsmath, amsfonts, amssymb, amsthm}
\usepackage{geometry}
\geometry{top=2cm, bottom=2cm, left=2.5cm, right=1.5cm}
\usepackage{hyperref}

\newtheorem{theorem}{Теорема}
\newtheorem{lemma}{Лемма}
\newtheorem{proposition}{Утверждение}
\renewcommand{\proofname}{Доказательство}

\title{Самоорганизующиеся структуры данных}
\date{}

\begin{document}

\maketitle

\section{Введение}
В данном конспекте рассматриваются самоорганизующиеся структуры данных (то есть те, которые могут менять свою структуру по некому правилу/алгоритму), а именно unsorted linear lists и их применение для задачи поддержания множества элементов. Пусть у нас есть список из элементов, где каждый элемент содержит указатель на следующий и предыдущий. Чтобы найти элемент, необходимо сделать линейный проход по списку, начиная с головы. Множество элементов списка обозначим $S$. В процессе обработки запросов алгоритм может динамически изменять порядок элементов в списке, чтобы оптимизировать время будущих обращений.

\section{Постановка задачи}
Рассмотрим последовательность запросов. Возможны следующие операции:
\begin{itemize}
    \item \textbf{Access(x).} Поиск элемента $x \in S$. Если $x$ находится на $k$-й позиции в списке, алгоритм проходит $k$ элементов, поэтому затраты на запрос равны $k$.
    \item \textbf{Insert(x).} Вставка элемента $x$. Алгоритм сканирует весь список (убеждаясь, что $x$ нет), а затем вставляет $x$ в конец. Затраты равны $n+1$, где $n$ --- текущая длина списка.
    \item \textbf{Delete(x).} Удаление элемента $x$. Затраты равны $k$.
\end{itemize}

Не умаляя общности, далее в анализе мы будем считать, что последовательность запросов состоит только из запросов типа \textbf{Access}. Все утверждения легко обобщаются и на \textbf{Insert}, и на \textbf{Delete}. 

\textbf{Модель стоимостей.} Сразу после доступа к элементу $x$, алгоритм может бесплатно переместить его на любую позицию ближе к началу списка (\textit{free exchanges}). Любой другой обмен двух соседних элементов в списке стоит 1 (\textit{paid exchanges}).

Мы оцениваем эффективность нашего онлайн-алгоритма $A$, сравнивая его с оптимальным офлайн-алгоритмом $OPT$, который заранее знает всю последовательность запросов $\sigma$. Пусть $C_A(\sigma)$ --- суммарные затраты алгоритма $A$, а $C_{OPT}(\sigma)$ --- затраты оптимального алгоритма. Теперь было бы неплохо для любого онлайн алгоритма придумать число, позволяющее сравнивать качество алгоритмов.

\textbf{Определение:} Детерминированный онлайн-алгоритм называется $c$-конкурентным ($c$-competitive), если существует константа $a$, такая что для любой последовательности запросов $\sigma$ выполняется:
$$C_A(\sigma) \le c \cdot C_{OPT}(\sigma) + a$$

\section{Детерминированные алгоритмы}
Рассмотрим три классических детерминированных онлайн-алгоритма:
\begin{itemize}
    \item \textbf{Move-to-front (MTF):} после обращения к элементу он всегда перемещается в самое начало списка.
    \item \textbf{Transpose:} запрошенный элемент меняется местами с непосредственно предшествующим ему элементом.
    \item \textbf{Frequency-count:} ведется подсчет частоты обращений к каждому элементу. Список всегда поддерживается отсортированным по невозрастанию частоты.
\end{itemize}

\begin{proposition}
Алгоритмы Transpose и Frequency-Count не являются $c$-конкурентными ни для какой константы $c$ (без доказательства).
\end{proposition}

\begin{proposition}
Для любого детерминированного $c$-конкурентного алгоритма $c \geq 2$ (без доказательства).
\end{proposition}

\begin{theorem}
Алгоритм Move-to-front является 2-конкурентным.
\end{theorem}

\begin{proof}
Рассмотрим последовательность запросов $\sigma = \sigma(1)\sigma(2)\dots\sigma(m)$. На каждом шаге $t$ будем сравнивать списки, полученные алгоритмами MTF и OPT. 
Будем называть \textit{инверсией} пару элементов $x, y$, расположенных в разном порядке в списках MTF и OPT (например, в OPT $x$ идет перед $y$, а в MTF --- наоборот).
Введем функцию потенциала $\Phi(t)$, равную количеству инверсий после обслуживания запроса $t$.

Пусть в момент $t$ запрошен элемент $x$. Пусть $k$ --- число элементов, стоящих строго перед $x$ в \textbf{обоих} списках, а $l$ --- число элементов, предшествующих $x$ в списке MTF, но находящихся после $x$ в списке OPT. 

Тогда фактические затраты алгоритмов на этом шаге:
$$C_{MTF}(t) = k + l + 1, \quad C_{OPT}(t) \ge k + 1$$

Оценим амортизированную стоимость шага для MTF: $a(t) = C_{MTF}(t) + \Phi(t) - \Phi(t-1)$.
Переместив $x$ в начало списка, MTF уничтожит ровно $l$ инверсий (с элементами, которые в OPT стоят после $x$) и создаст не более $k$ новых инверсий. Следовательно, изменение потенциала за счет действий MTF: $\Delta\Phi \le k - l$. 
Заметим также, что любой платный обмен (paid exchange), сделанный алгоритмом OPT, стоит ему 1, но может увеличить $\Phi$ не более чем на 1, что сохраняет итоговое неравенство.

Тогда:
$$C_{MTF}(t) + \Phi(t) - \Phi(t - 1) \le (k + l + 1) + (k - l) = 2k + 1 \le 2C_{OPT}(t) - 1$$

Суммируя это неравенство по всем $t$ от $1$ до $m$, и учитывая, что $\Phi(m) \ge 0$ и $\Phi(0) = 0$ (начальные списки совпадают), получаем $C_{MTF}(\sigma) \le 2C_{OPT}(\sigma)$, что доказывает 2-конкурентность.
\end{proof}

\section{Вероятностные алгоритмы}
\textbf{Определение:} Вероятностный онлайн алгоритм $A$ называется $c$-конкурентным, если существует константа $a$, что для любой последовательности операций $\sigma$:
$$\mathbb{E}[C_{A}(\sigma)] \le c \cdot C_{OPT}(\sigma) + a.$$

По сути то же определение, что и раньше, но мы усредняем стоимость по всем возможным вариантам действий алгоритма на заданной последовательности операций.

\begin{proposition}
Для любого вероятностного $c$-конкурентного алгоритма $c \geq 1.5$ (без доказательства).
\end{proposition}

\textbf{Алгоритм Timestamp(p)}. Пусть $p \in [0, 1]$ --- параметр алгоритма Timestamp (иногда будем сокращать до $TS$). Запрос $\sigma(t) = x$ обрабатывается следующим образом:

\begin{enumerate}
    \item Если элемент $x$ ни разу не запрашивался на интервале $[1, t-1]$, то его позиция в списке не меняется, и алгоритм завершает шаг.
    \item Иначе, пусть $t' \in [1, t-1]$ --- момент \textbf{предыдущего} запроса к $x$.
    \item С вероятностью $p$ выполняется \textbf{Шаг (a)}: элемент $x$ перемещается в самое начало списка.
    \item С вероятностью $1-p$ выполняется \textbf{Шаг (b)}:
    Алгоритм ищет элемент $v_x(t)$, который является \textbf{ближайшим к началу списка} среди всех элементов, \textbf{предшествующих} $x$ в текущий момент, и удовлетворяет хотя бы одному из условий:
    \begin{enumerate}
        \item[i)] он вообще не запрашивался в интервале $[t', t-1]$;
        \item[ii)] он запрашивался ровно один раз в интервале $[t', t-1]$, и этот запрос был обработан Шагом (b).
    \end{enumerate}
    Если такой элемент $v_x(t)$ найден, алгоритм вставляет $x$ непосредственно \textbf{перед} $v_x(t)$. Если такого элемента нет, мы считаем $v_x(t) = x$, то есть элемент $x$ просто остается на своем месте.
\end{enumerate}

\textit{Примечание: \textbf{Timestamp(p)} вынужден использовать информацию о прошлых запросах.}

Заметим, что при $p=1$ алгоритм вырождается в Move-to-front. 

\begin{theorem}
Алгоритм Timestamp(p) является $c$-конкурентным, где $c = \max(2 - p, 1 + p(2 - p))$.
\end{theorem}
\textit{Примечание:} При выборе $p = (3-\sqrt{5})/2$ достигается оптимальное для Timestamp(p) конкурентное отношение, равное золотому сечению $\phi \approx 1.62$, что лучше любого детерминированного алгоритма.

Для доказательства сформулируем вспомогательные леммы.

\begin{lemma}
Пусть $x$ и $y$ ($x \neq y$) --- элементы списка. Если $x$ запрашивается в момент $t'$, а затем в момент $t > t'$, и $y$ не запрашивается в интервале $[t', t]$, то сразу после обслуживания запроса $t$, элемент $x$ предшествует $y$ в списке.
\end{lemma}
\begin{proof}
Если $\sigma(t)$ обслуживается шагом (a), $x$ идет в начало списка и очевидно предшествует $y$. Если применяется шаг (b), $x$ вставляется перед $v_x(t)$, а значит, перед всеми элементами, которые не запрашивались на интервале $[t', t]$. Элемент $y$ удовлетворяет этому условию. Поскольку позиция $y$ не может поменяться (пока к нему не обратятся), $x$ будет предшествовать $y$.
\end{proof}

\begin{lemma}
Пусть в момент $t$ запрошен $x$, и $t' < t$ --- момент предыдущего запроса к $x$. Если элемент $y$ ($y \neq x$) удовлетворяет хотя бы одному из условий:
\begin{enumerate}
    \item[(a)] запрашивался хотя бы дважды на интервале $[t', t-1]$;
    \item[(b)] запрашивался ровно один раз на интервале $[t', t-1]$, и этот запрос был обработан Шагом (a);
\end{enumerate}
то $y$ предшествует $v_x(t)$ в списке в момент $t$.
\end{lemma}
\begin{proof}
Докажем пункт (a) от противного. Пусть $t_0$ --- самый первый момент, когда утверждение нарушается, $t'_0$ --- предыдущий запрос к $x$, а $z = v_x(t_0)$. Предположим, что $y$ запрашивался дважды, но не предшествует $z$. Пусть $t_y$ --- момент последнего запроса к $y$ на $[t'_0, t_0-1]$. 
Если $\sigma(t_y)$ обработан шагом (a), $y$ попадает в начало и точно предшествует $z$. Если шагом (b), вспомним, что по определению $v_x(t)$, элемент $z$ запрашивался максимум 1 раз (и через шаг b). Значит $z$ не может предшествовать $v_y(t_y)$. Тогда $y$ вставляется перед $z$ в момент $t_y$. 
Чтобы к моменту $t_0$ элемент $z$ снова оказался перед $y$, он должен быть запрошен в какой-то момент $t_z \in [t_y+1, t_0-1]$ и вставлен перед $y$. Но к моменту $t_z$ элемент $y$ уже запрашивался дважды с момента прошлого запроса к $z$, что противоречит выбору $t_0$ как самого первого момента нарушения леммы.

Случай (b) доказывается аналогично.
\end{proof}

\section{Схема доказательства Теоремы 2}
Введем индикаторную величину $C_{TS}(t, x) = 1$, если в момент $t$ элемент $x$ предшествует запрошенному элементу $\sigma(t)$, и $0$ иначе. Затраты алгоритма на всей последовательности можно расписать как:
$$C_{TS}(\sigma) = \sum_{t=1}^m \sum_{x \in S} C_{TS}(t, x)$$
Перегруппируем слагаемые, собирая их не по времени, а по неупорядоченным парам элементов $\{x, y\}$:
$$\mathbb{E}[C_{TS}(\sigma)] = \mathbb{E}\left[ \sum_{\{x,y\}} \left( \sum_{t:\sigma(t)=x} C_{TS}(t, y) + \sum_{t:\sigma(t)=y} C_{TS}(t, x) \right) \right]$$

Пусть $\sigma_{\{x,y\}}$ --- проекция последовательности запросов только на элементы $x$ и $y$ (то есть запросы берем из $\sigma$ только к $x$ и $y$). Пусть $\bar{C}_{OPT}(\sigma_{\{x,y\}})$ --- стоимость работы оптимального алгоритма на списке, состоящем только из этих двух элементов. Если расписать аналогично формуле выше, то легко будет увидеть, что:
$$C_{OPT}(\sigma) \ge \sum_{\{x,y\}} \bar{C}_{OPT}(\sigma_{\{x,y\}})$$

Такое разбиение позволяет рассматривать относительный порядок только двух элементов $\{x, y\}$, полностью игнорируя остальные элементы. Нам достаточно доказать, что для любой пары суммарное матожидание затрат Timestamp на этой паре не превосходит $c \cdot \bar{C}_{OPT}(\sigma_{\{x,y\}})$.

Последовательность $\sigma_{\{x,y\}}$ разбивается на фазы $P_j$. Фаза заканчивается, когда после серии запросов к одному элементу (хотя бы 2 запроса к одному и тому же) происходит запрос к другому. Фазы классифицируются на 4 типа (где $h \ge 2$):
\begin{enumerate}
    \item $P_j = x^h$ (или $y^h$)
    \item $P_j = x y^h$ (или $y x^h$)
    \item $P_j = (xy)^{h_1} x^{h_2}$ (или $(yx)^{h_1} y^{h_2}$), $h_1 \ge 1$
    \item $P_j = (xy)^{h_1} y^{h_2}$ (или $(yx)^{h_1} x^{h_2}$), $h_1 \ge 1$
\end{enumerate}

Используя доказанные Леммы 1 и 2, и вычисляя вероятности перемещений элементов на каждом шаге, можно оценить матожидание затрат Timestamp по сравнению с $\bar{C}_{OPT}$ для каждой фазы в отдельности (без доказательства, здесь на самом деле просто технические выкладки):
\begin{itemize}
    \item Для типа 1: $\mathbb{E}[C_{TS}(P_j)] \le (2 - p) \cdot \bar{C}_{OPT}(P_j)$
    \item Для типа 2: $\mathbb{E}[C_{TS}(P_j)] \le (1 + p(2 - p)) \cdot \bar{C}_{OPT}(P_j)$
    \item Для типов 3 и 4: $\mathbb{E}[C_{TS}(P_j)] \le \max(2 - p, 1 + p(2 - p)) \cdot \bar{C}_{OPT}(P_j)$
\end{itemize}

Поскольку любая фаза для любой пары элементов ограничивается коэффициентом $c = \max(2 - p, 1 + p(2 - p))$, суммирование этих оценок дает искомую конкурентность для всего алгоритма. Теорема доказана.

\end{document}